\section{Minor Test 2, 2016--17 (Semester 1)}

\subsection*{Question 1}
\textbf{Question.}
First, give a solution (algorithm) which allows at most \(n-1\) entries in
the buffer, of the producer--consumer problem with a finite buffer.
Secondly, modify the algorithm to remedy this deficiency.

\textbf{Solution.}
Let \texttt{buffer[0..n-1]} be a circular buffer.  Initially,
\(\texttt{in}=\texttt{out}=0\).  For one producer and one consumer, the
following algorithm reserves one array position in order to distinguish
``empty'' from ``full'':

\begin{verbatim}
/* Producer */
for (;;) {
    item = produce_item();
    while ((in + 1) % n == out)
        ;                       /* full: spin */
    buffer[in] = item;
    in = (in + 1) % n;
}

/* Consumer */
for (;;) {
    while (in == out)
        ;                       /* empty: spin */
    item = buffer[out];
    out = (out + 1) % n;
    consume_item(item);
}
\end{verbatim}

Here \(\texttt{in}=\texttt{out}\) denotes an empty buffer, whereas
\((\texttt{in}+1)\bmod n=\texttt{out}\) denotes a full buffer.  Thus at
most \(n-1\) items can be stored: one physical slot is deliberately
wasted.  (With multiple producers or multiple consumers, the index
updates would additionally need mutual exclusion.)

To use all \(n\) slots, keep empty and full states separately using
counting semaphores.  Let
\[
  \texttt{mutex}=1,\qquad \texttt{empty}=n,\qquad \texttt{full}=0,
  \qquad \texttt{in}=\texttt{out}=0.
\]
The corrected bounded-buffer algorithm is:

\begin{verbatim}
/* Producer */
for (;;) {
    item = produce_item();
    sem_wait(&empty);
    sem_wait(&mutex);
    buffer[in] = item;
    in = (in + 1) % n;
    sem_post(&mutex);
    sem_post(&full);
}

/* Consumer */
for (;;) {
    sem_wait(&full);
    sem_wait(&mutex);
    item = buffer[out];
    out = (out + 1) % n;
    sem_post(&mutex);
    sem_post(&empty);
    consume_item(item);
}
\end{verbatim}

\texttt{empty} now records the number of free slots and \texttt{full}
records the number of occupied slots, so the state with
\(\texttt{full}=n\) is unambiguous.  Hence all \(n\) slots can be used.
\texttt{mutex} protects the buffer and indices, so this version also
works with multiple producers and consumers.

\subsection*{Question 2}
\textbf{Question.}
Apply the deadlock-detection algorithm to
\[
 \mathit{Available}=(2,1,0,0),
\]
\[
 \mathit{Request}=
 \begin{bmatrix}
 2&0&0&1\\
 1&0&1&0\\
 2&1&0&0
 \end{bmatrix},
 \qquad
 \mathit{Allocation}=
 \begin{bmatrix}
 0&0&1&0\\
 2&0&0&1\\
 0&1&2&0
 \end{bmatrix},
\]
and show the result.

\textbf{Solution.}
Call the three rows \(P_1,P_2,P_3\).  Initially,
\[
 \mathit{Work}=(2,1,0,0).
\]
Every process has a nonzero allocation, so initially
\(\mathit{Finish}[1]=\mathit{Finish}[2]=\mathit{Finish}[3]=\mathrm{false}\).
A vector comparison below is componentwise.

In the first scan,
\[
\begin{aligned}
 \mathit{Request}_1=(2,0,0,1)&\nleq(2,1,0,0)
       &&\text{(the fourth component fails)},\\
 \mathit{Request}_2=(1,0,1,0)&\nleq(2,1,0,0)
       &&\text{(the third component fails)},\\
 \mathit{Request}_3=(2,1,0,0)&\leq(2,1,0,0).
\end{aligned}
\]
Therefore \(P_3\) can finish and release its allocation:
\[
 \mathit{Work}
 =(2,1,0,0)+(0,1,2,0)
 =(2,2,2,0),
 \qquad \mathit{Finish}[3]=\mathrm{true}.
\]

With this new work vector, \(P_1\) still cannot finish because it needs
one unit of the fourth resource.  However,
\[
 \mathit{Request}_2=(1,0,1,0)\leq(2,2,2,0),
\]
so \(P_2\) finishes and releases its allocation:
\[
 \mathit{Work}
 =(2,2,2,0)+(2,0,0,1)
 =(4,2,2,1),
 \qquad \mathit{Finish}[2]=\mathrm{true}.
\]
Finally,
\[
 \mathit{Request}_1=(2,0,0,1)\leq(4,2,2,1),
\]
so \(P_1\) finishes:
\[
 \mathit{Work}
 =(4,2,2,1)+(0,0,1,0)
 =(4,2,3,1),
 \qquad \mathit{Finish}[1]=\mathrm{true}.
\]
Thus a complete finishing sequence is
\[
 \boxed{P_3,\ P_2,\ P_1}.
\]
All \(\mathit{Finish}\) entries become true; therefore no process is
deadlocked.

\subsection*{Question 3}
\textbf{Question.}
Suppose there are two types of philosophers.  A ``lefty'' always picks
up the left fork first, while a ``righty'' always picks up the right
fork first.  Using semaphores, define both types of philosophers'
behaviour.

\textbf{Solution.}
Number the \(N\) philosophers and forks from \(0\) to \(N-1\).  For
philosopher \(i\), define
\[
  \mathit{left}(i)=i,\qquad
  \mathit{right}(i)=(i+1)\bmod N.
\]
Each \texttt{fork[i]} is a binary semaphore initialized to \(1\).

\begin{verbatim}
void lefty(int i)
{
    int left  = i;
    int right = (i + 1) % N;

    for (;;) {
        think();
        sem_wait(&fork[left]);
        sem_wait(&fork[right]);
        eat();
        sem_post(&fork[right]);
        sem_post(&fork[left]);
    }
}

void righty(int i)
{
    int left  = i;
    int right = (i + 1) % N;

    for (;;) {
        think();
        sem_wait(&fork[right]);
        sem_wait(&fork[left]);
        eat();
        sem_post(&fork[left]);
        sem_post(&fork[right]);
    }
}
\end{verbatim}

If all philosophers acquire forks in the same order, each can hold one
fork while waiting for the next, producing a circular wait.  If there
is at least one lefty and at least one righty (for example, make the
last philosopher right-handed and all others left-handed), this
directed cycle is broken, so deadlock is avoided.  This argument
establishes deadlock freedom; it does not by itself guarantee freedom
from starvation under an unfair semaphore implementation.

\subsection*{Question 4}
\textbf{Question.}
Five batch jobs \(A\) through \(E\) arrive at essentially the same time.
Their running times are \(15,9,3,6,12\) minutes, respectively, and
their priorities are \(6,3,7,9,4\), respectively, where a lower value
means a higher priority.  For non-preemptive priority scheduling and
shortest-job-first (SJF), determine the turnaround time of each process
and the average turnaround time.

\textbf{Solution.}
All arrival times are \(0\), so for every job,
\[
  \mathit{turnaround\ time}
  =\mathit{completion\ time}-\mathit{arrival\ time}
  =\mathit{completion\ time}.
\]

\paragraph{Non-preemptive priority scheduling.}
Increasing priority number gives the execution order
\[
 \underbrace{0}_{}
 \;|\;B\;|\;9\;|\;E\;|\;21\;|\;A\;|\;36\;|\;C\;|\;39\;|\;D\;|\;45.
\]
Thus
\[
\begin{array}{c|ccccc}
\text{Job} & A&B&C&D&E\\ \hline
\text{Turnaround time (min)}&36&9&39&45&21
\end{array}
\]
and
\[
 \overline{T}_{\rm priority}
 =\frac{36+9+39+45+21}{5}
 =\frac{150}{5}
 =\boxed{30\text{ min}}.
\]

\paragraph{Non-preemptive SJF.}
Increasing burst time gives the execution order
\[
 \underbrace{0}_{}
 \;|\;C\;|\;3\;|\;D\;|\;9\;|\;B\;|\;18\;|\;E\;|\;30\;|\;A\;|\;45.
\]
Thus
\[
\begin{array}{c|ccccc}
\text{Job} & A&B&C&D&E\\ \hline
\text{Turnaround time (min)}&45&18&3&9&30
\end{array}
\]
and
\[
 \overline{T}_{\rm SJF}
 =\frac{45+18+3+9+30}{5}
 =\frac{105}{5}
 =\boxed{21\text{ min}}.
\]

\subsection*{Question 5}
\textbf{Question.}
What are the four conditions that create deadlock?  List two ways in
which the no-preemption condition can be prevented.

\textbf{Solution.}
The four Coffman conditions, all of which must hold simultaneously for
deadlock to be possible, are:
\begin{enumerate}
 \item \textbf{Mutual exclusion:} at least one resource is
       non-shareable, so only one process can use it at a time.
 \item \textbf{Hold and wait:} a process holds one or more resources
       while waiting to acquire additional resources.
 \item \textbf{No preemption:} a resource cannot be forcibly taken
       away; it is released only voluntarily by its holder.
 \item \textbf{Circular wait:} there is a cycle of processes in which
       each process waits for a resource held by the next process in
       the cycle.
\end{enumerate}

Deadlock can be prevented by invalidating the no-preemption condition,
for example in either of the following ways.
\begin{enumerate}
 \item If a process holding resources requests another resource that
       cannot be allocated immediately, require it to release all
       resources it currently holds.  It retries only after the needed
       resources can be acquired again.  Thus it cannot indefinitely
       retain its allocation while blocked.
 \item Permit the operating system to forcibly reclaim a resource from
       its current holder, return that resource to the available pool,
       and later restart or roll back the affected computation.  This
       is practical only for resources whose state can be saved and
       restored, such as CPU state and, with appropriate mechanisms,
       memory pages; it is generally unsuitable for inherently
       non-preemptible devices in the middle of an operation.
\end{enumerate}
The rule that a process must request all resources at once is instead
normally classified as prevention of \emph{hold and wait}, not as
prevention of no-preemption.
